\section{Betti and De Rham moduli spaces}\label{section: de rham moduli space}

We now come to the study of the Betti and de Rham moduli spaces of the stacky curve $\mc{X}$.
Recall that
\begin{gather}
\mc{M}_{n, DR}(\mc{X}) \coloneqq \left< (\mc{E}, \nabla) \mymid \mc{E}\in \Bun_n(\mc{X}), \; \nabla\colon \mc{E} \to \mc{E}\otimes \Omega_{X} \text{ is a flat connection} \right> \\
\mc{M}_{n, B}(\mc{X}) \coloneqq \text{Rep}(\pi_1(\mc{X}, x), GL_n).
\end{gather}
The Riemann-Hilbert correspondence establishes an isomorphism
% natural isomorphism
% \begin{equation}
% \mc{M}_{DR}(\mc{X}) \cong \mc{M}_B(\mc{X}),
% \end{equation}
% inducing an isomorphism
of the coarse moduli spaces
\begin{equation} \label{eq: riemann holbert correspondence}
M_{n, DR}(\mc{X})^{\text{an}} \cong M_{n, B}(\mc{X})^{\text{an}}.
\end{equation}
If we fix a rank $n$ and use the presentation of $\pi_1(\mc{X}, x)$ found in \eqref{eq: fundamental group of stacky curve}, we get
\begin{equation} \label{eq: de rham moduli space rnk n}
M_{n, B}(\mc{X}) \cong \left\{ ((A_i, B_i)_{i = 1}^g, C) \in GL_n^{2g +1} \mymid \prod_{i=1}^{g}[A_i, B_i] = C, \, C^n = Id \right\} // GL_n.
\end{equation}
As already observed for the moduli space of vector bundles $\Bun(\mc{X})$ and for the Dolbeault moduli space, also the de Rham moduli space admits a decomposition in connected components
\begin{equation}
\mc{M}_{n, DR}(\mc{X}) = \coprod_{\underline{d}} \mc{M}_{n, DR}^{\underline{d}}(\mc{X}) = \coprod_{(\underline{m}, d)} \mc{M}_{\underline{m}, DR}^d(\mc{X}),
\end{equation}
with the latter union ranging over the pairs $(\underline{m}, d)$ such that $\sum_{i=0}^{n-1}m_{i} = n$. 
Recall that for a vector bundle to admit a flat connection it is a necessary condition that $\deg \mc{E} = 0$. Therefore we can set $\mc{M}_{\underline{m}, DR} \coloneqq \mc{M}_{\underline{m}, DR}^0(\mc{X})$ and get
\begin{equation}
\mc{M}_{n, DR}(\mc{X}) = \coprod_{\underline{m}} \mc{M}_{\underline{m}, DR}(\mc{X}).
\end{equation}

This decomposition gives decomposition in connected components of the coarse moduli space
\begin{equation} \label{eq: de rham moduli space, decomposition}
M_{n, DR}(\mc{X}) = \coprod_{\underline{m}} M_{\underline{m}, DR}(\mc{X}).
\end{equation}
We now want to understand how this decomposition of the coarse moduli space $M_{n, DR}$ translates to the Betti moduli space, and in particular to the character variety on the right hand side of \eqref{eq: de rham moduli space rnk n}.

\subsection{The monodromy representation}
The isomorphism in \zcref{eq: riemann holbert correspondence} is obtained by mapping a pair $(\mc{E}, \nabla)$ to the associated monodromy representation at $x$. In particular, the matrices $A_i, \, B_i, \, C$ in \eqref{eq: de rham moduli space rnk n} are obtained by the action of the generators $a_i, \, b_i, \, c$ of the fundamental group of $\mc{X}$, as expressed in \zcref{stacky curve}, and are well defined up to conjugation.

The multiplicity vector $\underline{m}(\mc{E})$ depends only on the behavior of $\mc{E}$ at the stacky point, and therefore can be computed in a neighborhood of $p$ isomorphic to $[\Ab^1 / \mu_n]$. In particular, looking back at the Seifert-Van Kampen computation of $\pi_1(\mc{X}, x)$ in \zcref{stacky curve}, we see that the decomposition \eqref{eq: de rham moduli space, decomposition} will translate on the character variety only in terms of conditions on the matrix $C$.

In particular, after restricting to $[\Ab^1/ \mu_n]$, we can find the matrix $C$ as obtained by the monodromy action of the generator of $\pi_1([\Ab^1 / \mu_n], x) \cong \mu_n$ on $(\mc{E}\vert_{[\Ab^1 / \mu_n]}, \nabla\vert_{[\Ab^1 / \mu_n]})$.

Therefore, let us consider the case of a vector bundle $\mc{E}$ on $[\Ab^1 / \mu_n]$, with a flat connection $\nabla$.
As mentioned in \zcref{stacky curve}, the projection map
\begin{equation}
p \colon \Ab^1 \to [\Ab^1 / \mu_n]
\end{equation}
is a principal $\mu_n$-bundle and it is the universal cover of $[\Ab^1 / \mu_n]$. The datum of the pair $(\mc{E}, \nabla)$ is equivalent to the $\mu_n$-equivariant pair $(p^*\mc{E}, p^*\nabla)$ on $\Ab^1$. The fundamental group of $[\Ab^1 / \mu_n]$ is identified with $\mu_n$ and under this identification the monodromy action on $(\mc{E}, \nabla)$ is the pushforward of the $\mu_n$-action on $(p^*\mc{E}, p^*\nabla)$.

Since $\Ab^1$ is simply connected and $p^*\nabla$ is a flat connection on $p^*\mc{E}$, the pair $(p^*\mc{E}, p^*\nabla)$ is actually equivalent to the trivial bundle, with the standard derivation. Moreover, since $\Ab^1$ is affine, this is just the sheafification of the pair $(\Cb[x]^n, \dd)$. The $\mu_n$-equivariant structure corresponds to an action
\begin{equation}
\mu_n \to \Aut_{\Cb[x]}( (\Cb[x]^n, \dd) ).
\end{equation}
Notice that an automorphism $\varphi\colon \Cb[x]^n \to \Cb[x]^n$ fixes the differential $\dd$ if and only if $\dd(\varphi) = 0$. Therefore we have the following isomorphisms
\begin{equation}
\begin{tikzcd}
	{\Aut_{\Cb[x]}(\Cb[x]^n)} & {GL_n(\Cb[x])} \\
	{\Aut_{\Cb[x]}((\Cb[x]^n, \dd))} & {GL_n(\Cb).}
	\arrow["\cong"{marking, allow upside down}, draw=none, from=1-1, to=1-2]
	\arrow["\subseteq"{marking, allow upside down}, draw=none, from=2-1, to=1-1]
	\arrow["\cong"{marking, allow upside down}, draw=none, from=2-1, to=2-2]
	\arrow["\subseteq"{marking, allow upside down}, draw=none, from=2-2, to=1-2]
\end{tikzcd}
\end{equation}
In particular, the $\mu_n$-equivariant structure on $(p^*\mc{E}, p^*\nabla)$ gives an action
\begin{equation}
\rho \colon \mu_n \to GL_n(\Cb),
\end{equation}
i.e. an $n$-dimensional representation of $\mu_n$, which is exactly the monodromy representation of the pair $(\mc{E}, \nabla)$.

We now want to compare this with the representation obtained as $\mc{E}\vert_{B\mu_n}$. Consider
\begin{equation}
\begin{tikzcd}
	{* = \{ 0 \}} & {\Ab^1} \\
	{B\mu_n} & {[\Ab^1 / \mu_n]}
	\arrow["i", hook, from=1-1, to=1-2]
	\arrow[from=1-1, to=2-1]
	\arrow["p", from=1-2, to=2-2]
	\arrow["j"', hook, from=2-1, to=2-2]
\end{tikzcd}\
\end{equation}
The representation associated to $\mc{E}\vert_{B\mu_n}$ is obtained by the $\mu_n$-equivariant structure of $(p\vert_{\{0\}})^* \mc{E}\vert_{B\mu_n}$, corresponding to a morphism
\begin{equation}
\mu_n \to \Aut_{\{0\}}((p\vert_{\{0\}})^* \mc{E}\vert_{B\mu_n}) = \Aut_{\{0\}}( i^* p^* \mc{E} ).
\end{equation}
Under the identifications $\qcoh(\Ab^1) \cong \Cb[x]\mh\Mod$ and $\qcoh(*) \cong \Cb\mh\Mod$, the restriction $i^*$ is just the morphism of evaluation at $0$. Since $p^* \mc{E}$ is the trivial rank $n$ vector bundle, we get
\begin{equation}
\begin{tikzcd}
	{\mu_n} & {GL_n(\Cb[x])} & {\Aut_{\Ab^1}(p^* \mc{E})} \\
	& {GL_n(\Cb)} & {\Aut_{\{0\}}(i^* p^* \mc{E}).}
	\arrow["{\rho'}", from=1-1, to=1-2]
	\arrow["\sigma"', from=1-1, to=2-2]
	\arrow["\cong"{marking, allow upside down}, draw=none, from=1-2, to=1-3]
	\arrow["{{\text{ev}_0}}", from=1-2, to=2-2]
	\arrow["{{i^*}}", from=1-3, to=2-3]
	\arrow["\cong"{marking, allow upside down}, draw=none, from=2-2, to=2-3]
\end{tikzcd}
\end{equation}
But $\rho'$ must factor through the inclusion
\begin{equation}
\Aut_{\Ab^1}( (p^*\mc{E}, \nabla) ) \subseteq \Aut_{\Ab^1}(p^*\mc{E}),
\end{equation}
so we get
\begin{equation}
\begin{tikzcd}
	{\mu_n} & {GL_n(\Cb)} & {GL_n(\Cb[x])} \\
	&& {GL_n(\Cb).}
	\arrow["\rho", from=1-1, to=1-2]
	\arrow["{\rho'}", curve={height=-18pt}, from=1-1, to=1-3]
	\arrow["\sigma"', from=1-1, to=2-3]
	\arrow["\subseteq"{marking, allow upside down}, draw=none, from=1-2, to=1-3]
	\arrow["{{\text{ev}_0}}", from=1-3, to=2-3]
\end{tikzcd}
\end{equation}
Clearly the composition $GL_n(\Cb) \subseteq GL_n(\Cb[x]) \xrightarrow{\text{ev}_0} GL_n(\Cb)$ is the identity, so the representations $\rho$ and $\sigma$ coincide.

In conclusion, we obtained the following result.
\begin{prop} \label{prop: monodromy representation}
The monodromy representation of a pair $(\mc{E}, \nabla)$ coincides with the representation obtained as $\mc{E}\vert_{B \mu_n}$.
\end{prop}

\subsection{Betti moduli space}

Applying \zcref{prop: monodromy representation} to the computation of the coarse moduli spaces, we get that
\begin{equation}
M_{\underline{m}, B}(\mc{X}) \cong \left\{ (A_i, B_i)_{i = 1}^g \in GL_n^{2g} \mymid \prod_{i=1}^{g}[A_i, B_i] = C_{\underline{m}} \right\} // GL_n,
\end{equation}
where $C_{\underline{m}}$ is a matrix representing the action of the generator of $\mu_n$ in the representation defined by the multiplicity vector $\underline{m}$.
In particular, $C_{\underline{m}}$ can be taken as a diagonal matrix with exactly $m_l$ entries equal to $\mathrm{e}^{\frac{2 \pi \mathrm{i}}{n} l}$.

In the case $\underline{m} = n_{-d}$, we get

\begin{prop}
The $n_{-d}$ stratum of coarse moduli space of the Betti moduli stack of $\mc{X}$ is isomorphic to the following variety
\begin{equation}
M_{n_{-d}, B}(\mc{X}) \cong \left\{ (A_i, B_i)_{i = 1}^g \in GL_n^{2g} \mymid \prod_{i=1}^{g}[A_i, B_i] = \mathrm{e}^{-\frac{2 \pi \mathrm{i}}{n} d} Id \right\} // GL_n.
\end{equation}
\end{prop}